\subsection{動作原理}

学校レベルの数学から、9による除算は$\frac{1}{9}$による乗算に置き換えることができることを覚えています。
実際、コンパイラは浮動小数点演算でそうすることがあります。例えば、x86コードの\INS{FDIV}命令は\INS{FMUL}に置き換えることができます。
少なくともMSVC 6.0は9による除算を$0.111111...$による乗算に置き換えるため、元のソースコードにどの操作があったかを確信するのは困難な場合があります。

しかし、整数値と整数CPUレジスタを扱う場合、分数を使用することはできません。
ただし、分数を次のように書き直すことができます：

% FIXME: equation size
\begin{center}
$result = \frac{x}{9} = x \cdot \frac{1}{9} = x \cdot \frac{1 \cdot MagicNumber}{9 \cdot MagicNumber}$
\end{center}

$2^n$による除算が非常に高速（シフトを使用）であることを考慮すると、次の方程式が成り立つような$MagicNumber$を見つける必要があります：$2^n = 9 \cdot MagicNumber$。

$2^{32}$による除算は多少隠されています：EAXの積の下位32ビットは使用されず（破棄され）、積の上位32ビット（EDX内）のみが使用され、さらに1ビット追加でシフトされます。

つまり、先ほど見たアセンブリコードは{\Large $\frac{954437177}{2^{32+1}}$}で乗算、または{\Large $\frac{2^{32+1}}{954437177}$}で除算しています。
除数を求めるには、分子を分母で割るだけです。
Wolfram Alphaを使用すると、8.99999999....という結果（9に近い）を得ることができます。

% TODO ref to https://yurichev.com/blog/signed_division_using_shifts/

詳細については\InSqBrackets{\HenryWarren 10-3}を参照してください。

理解を深めるためのいくつかの言葉。
多くの人が「隠された」$2^{32}$または$2^{64}$による除算を見逃します。これは積の下位32ビット部分（または64ビット部分）が使用されない場合です。
また、ここで逆元の剰余が使用されているという誤解もあります。これは近いですが、同じではありません。
通常、\textit{マジック係数}を見つけるために拡張ユークリッド算法が使用されますが、実際には、この算法は方程式を解くために使用されます。他の方法を使用して解くこともできます。
いずれにせよ、拡張ユークリッド算法はおそらくそれを解く最も効率的な方法です。
また、言うまでもなく、一部の除数では方程式が解けません。これは結果が整数のみであることを許可するディオファンティン方程式であり、結局のところ整数CPUレジスタで作業するからです。